\documentclass[10pt]{article}
\usepackage[a4paper,margin=0.7in]{geometry}
\linespread{0.98}
\usepackage{parskip}
\usepackage[hidelinks]{hyperref}
\pagestyle{empty}

\begin{document}

\noindent
Quang-Vinh Dang\\
British University Vietnam\\
Hung Yen, Vietnam\\
\href{mailto:vinh.dq4@buv.edu.vn}{vinh.dq4@buv.edu.vn}

\vspace{0.7em}
\noindent 25 July 2026

\vspace{0.7em}
\noindent
The Editors\\
\emph{Computers \& Security}\\
Elsevier

\vspace{0.8em}
\noindent Dear Editors,

\vspace{0.6em}
I am pleased to submit my manuscript, \textbf{``Conformal Surprisal Fusion for
Detecting Prompt-Injection Compromise in LLM Agent Traces,''} for consideration as
a research article in \emph{Computers \& Security}.

Large language model agents that call external tools now read email, browse the
web and move money on a user's behalf. Because they ingest untrusted content and
fold it into the context that drives their next action, they are exposed to
indirect prompt injection. Runtime monitors that watch only the black-box execution
trace need no model access and no retraining, but the literature has converged
almost entirely on one kind of evidence: whether untrusted text \emph{looks like an
instruction}. This paper asks what else the trace affords, how heterogeneous
evidence should be combined, and where the information in a trace runs out.

The work makes four contributions. First, it introduces \textbf{SENTRY-Fuse}, a
monitor that converts heterogeneous trace signals---one scoring the observation
stream, two auditing the action stream---into conformal tail surprisals against a
benign calibration split and fuses them by addition. The construction is
scale-free, non-decreasing in every component and free of tuned signal weights,
and the paper proves that the natural alternative of subtracting the benign
training maximum is many-to-one and caps AUROC at one half. That is not
hypothetical: an earlier version of this work did exactly that, at a cost of $0.25$
AUROC and a confidently wrong explanation of the resulting failure. Second,
action-side evidence is shown to be \emph{complementary} rather than redundant:
over $90$ genuinely compromised trajectories the fusion reaches
$\mathrm{AUROC} = 0.957 \pm 0.004$ and a true-positive rate of $0.919 \pm 0.007$ at
a realised $3.4\%$ false-positive rate, against $0.944$ and $0.729$ for the
instruction signal alone, and it cuts that signal's across-split standard deviation
from $\pm 0.364$ to $\pm 0.007$. Third, it formalises the \textbf{detectability
ceiling} of trace-only monitoring and then measures it: on AgentDojo the fraction
of compromises leaving no trace evidence is \emph{zero}, because the benchmark
delivers payloads through tool outputs a successful injection must have read. The
ceiling is an artefact of benchmark construction rather than a limit of the
problem---a negative result about the field's evaluation practice, and the finding
I expect to be most useful to others. Fourth, it documents a label-polarity
pitfall in a widely used harness that inverts the compromised and resisted
classes, with the ground-truth procedure that resolves it.

I should be explicit that the paper does \emph{not} claim to beat the state of the
art. The strongest reported trace-only figure ($95.75\%$ true positives at
$3.66\%$ false positives) comes from a preprint version that later versions
withdraw, was measured on a different corpus with a different agent, and I
re-implemented no baseline on my own corpus. My numbers sit in the same regime,
and I say so rather than constructing a comparison the evidence does not support.

I believe the work suits \emph{Computers \& Security} in particular. It is a
security paper first: it defines a threat model, evaluates a defence against real
attacks on real agent traces, and reports its limits rather than only its
successes. It offers two things of direct use to others: a measurement showing that
the benchmarks the field relies on cannot exhibit the hardest case, and a
methodological warning that affects how results on a standard benchmark should be
read. Its negative results are reported in full---four signals I implemented and
rejected, two of which an earlier version of this work presented as part of the
method, and the finding that a hand-picked threshold outperformed the anytime-valid
machinery from which the project began.

All code, the trajectory corpora, the cached model-judge scores and the analysis
scripts that regenerate every number, table and figure are publicly released, and
the evaluation reproduces offline without access to any model API.

I confirm that this manuscript is original, has not been published previously and
is not under consideration elsewhere. There are no competing interests and no
external funding was received; as sole author I am responsible for all aspects of
the work. Thank you for your consideration, and I look forward to the reviewers'
comments.

\vspace{0.6em}
\noindent Sincerely,

\vspace{0.4em}
\noindent Quang-Vinh Dang

\end{document}
